\title{Summary of Revisions}
\vskip 0.2in

\noindent Dear IJRR Editors and Reviewers,
\vskip 0.1in

\noindent An earlier version of this manuscript received detailed feedback from three reviewers at IEEE Transactions on Robotics (T-RO Manuscript 25-1489), with the editor concluding that the manuscript at that time did not yet meet the bar for ``new results of substantive research significance and impact beyond the previous papers.'' We took the reviewers' technical concerns seriously and have made substantial efforts in this revision to strengthen the manuscript along every dimension they identified. This document is provided as supplementary material to make the revision history transparent and to help the IJRR review process by showing precisely how each prior comment was addressed.

\vskip 0.1in

\noindent At the structural level, the contribution narrative has been reworked so that the framework's complete technical scope is visible end-to-end. The conference-introduced scalability machinery (the high-level manager and symbolic repair with dynamic feasibility checks) is now presented alongside---and at the same level of detail as---the journal-specific mechanisms developed for this article: delay-aware coordination of the strategy automaton with the tracking controller, kinematic-feasibility re-targeting in the online MICP, and collision-region selection with swing-foot clearance constraints. Each of the journal-specific mechanisms is discussed in the introduction with explicit motivation for why it closes a concrete deployment gap, rather than presented as a standalone novelty bullet.

\vskip 0.1in

\noindent The major changes in this submission are:

\noindent \textbf{1. Contribution narrative.} The introduction now lists four contributions: (i) the integrated reactive-synthesis-and-MICP framework; (ii) scalable offline synthesis via the high-level manager and symbolic repair, with concrete $80.9$--$90.1\%$ state-space-reduction and $71.6$--$97.6\%$ MICP-solve-reduction numbers; (iii) the online-execution module with new mechanisms introduced in this article---kinematic-feasibility re-targeting, delay-aware coordination between the strategy automaton and the tracking controller, and collision-region selection in the online MICP; (iv) comprehensive empirical validation including pure-MIP and heuristic baselines, hardware deployment on two robot platforms, and a preliminary yaw-extension feasibility study. The introduction also includes a dedicated paragraph that motivates each of the journal-specific online-execution mechanisms with explicit discussion of the deployment gap it closes.

\noindent \textbf{2. Online-execution clarity.} Section~\ref{sec:online_execution} (online execution) was the most-criticized section in the prior round. The collision-avoidance subsection has been rewritten to explicitly distinguish between collision-free region selection and swing-foot clearance, name which robot parts are checked, and acknowledge that perception is out of scope. A new short subsection explains why symbolic--continuous interaction operates at the transition timescale rather than per-NMPC-tick.

\noindent \textbf{3. Quantitative presentation overhaul, with comprehensive results for symbolic repair.} Several reviewers asked for more quantitative evidence on how often and how expensively repair is performed; the IJRR submission addresses this at both the offline and runtime levels. Table~\ref{tab:offline_synthesis_time} (\textit{Offline Synthesis and Repair Time}) reports per-scenario success/total state-pair counts, original/new/total skill counts, symbolic-repair wall-clock times, and gait-fixed/gait-free feasibility-check times: across the four scenarios, repair successfully resolves $93.4\%$ of all terrain--request state pairs, and the number of newly discovered skills is reduced by $71.6$--$97.6\%$ relative to exhaustive enumeration. A new \textit{Runtime Repair Case Study} subsection (Sec.~\ref{sec:runtime_repair_case_study}) reports concrete timing for the two run-time scenarios---unforeseen terrain ($12.36$\,s total, $0.18$\,s symbolic repair plus $12.18$\,s gait-free MICP) and online MICP solving failure ($0.11$\,s and $0.07$\,s for the 7-type and 14-type rebar cases). The introduction's scalability-motivation paragraphs cite the manager's $80.9$--$90.1\%$ symbolic state-space reduction directly. On the implementation side, the results section now declares the solver (Gurobi~9), states the hardware platform (i9-13900K, 64\,GB RAM), defines the pure-MIP baseline precisely, and adds a std-dev column to the hardware MICP-solve-time table.

\noindent \textbf{4. Strengthened technical soundness of the newly proposed mechanisms via a new benchmark and explicit ablations.} The Pure-MIP benchmark subsection is substantially expanded. Fig.~\ref{fig:pure_mip_benchmark} now reports per-MIP solve time across four terrains and five planning horizons as box plots aggregated over $10$ seeded waypoints per scenario. A new end-to-end benchmark on the dense-rebar scenario (Table~\ref{tab:pure_mip_benchmark}) compares the proposed planner against pure-MIP $H{=}2$\,s and $H{=}3$\,s baselines and against two ablations of the newly proposed mechanisms (disabling the MIP feasibility check; disabling delay-aware coordination), with success rate, mean traversal time, and per-module computational cost reported per trial. The proposed planner matches the pure-MIP $H{=}3$\,s success rate ($90\%$) at $\sim 75\%$ of its traversal time and $<50\%$ of its MIP-transition cost; disabling the feasibility check drops the success rate from $90\%$ to $60\%$, and disabling delay-aware coordination inflates the mean traversal time from $18.5$\,s to $23.4$\,s---directly demonstrating that each newly proposed mechanism is necessary for the reported performance. A representative qualitative comparison is shown in Fig.~\ref{fig:dense_rebar_comparison}.

\noindent \textbf{5. Calibrated claims and honest limitations.} We tone down ``dynamically changing environments,'' ``dynamic locomotion,'' and ``safety guarantees'' wherever the experimental evidence does not warrant the strongest reading. We disclose the angular-momentum simplification of the MICP dynamics. The yaw extension is demoted to a preliminary feasibility study with explicit caveats. The discussion section is expanded with paragraphs on perception uncertainty, stronger future-baseline comparisons (Winkler 2018 phase-based parameterization and perceptive RL controllers), and bipedal-extension considerations.

\vskip 0.1in

\noindent All revisions to the manuscript are highlighted \textcolor{blue}{in blue text} to differentiate them from prior content. In the per-reviewer sections that follow, the original reviewer comments are boxed and our responses are detailed underneath, with revised passages quoted in blue italic for the reviewer's convenience. Unless otherwise noted, references to sections, figures, and equations in our responses point to the IJRR manuscript; references inside the boxed reviewer comments point to the prior T-RO version of the manuscript.

\vskip 0.2in

\noindent Sincerely, \\
\noindent Ziyi Zhou, Qian Meng, Hadas Kress-Gazit, and Ye Zhao
